
\subsection{式 (3.8.36) 的证明}
\label{subsec:7-3-3}

本小节证明定理~\ref{thm:3-8-32}
中的 (3.8.35) 和 (3.8.36)。
(3.8.35) 的证明，
是 (3.8.33)、(3.8.34)
及定理~\ref{thm:7-3-1} 证明的直接推广；
因此重点证明 (3.8.36)。

设 $Q_1,\ldots,Q_{\ell_1}$
为 $M$ 上的光滑奇异单形。
取 $M$ 的基本类，
仿照式~\eqref{eq:7-3-6}
以单纯剖分表示为
\begin{equation}
[M]=\sum_{a\in A}\pm\Delta_a^{2n}.
\tag{7.3.31}\label{eq:7-3-31}
\end{equation}
置
\begin{align}
\vec Q
&=Q_1\otimes\cdots\otimes Q_{\ell_1}
\in B_{\ell_1}C(M;\Lambda_{0,\mathrm{nov}}),
\tag{7.3.32.1}\label{eq:7-3-32-1}\\
\vec Q^+
&=\vec Q\otimes
\underbrace{M\otimes\cdots\otimes M}_{\ell_2\text{ 个}}
\in B_\ell C(M;\Lambda_{0,\mathrm{nov}}),
\tag{7.3.32.2}\label{eq:7-3-32-2}
\end{align}
其中 $\ell=\ell_1+\ell_2$。
\translatornote{源书在 (7.3.32.2)
把 $\vec Q^+$ 记为 $B_{\ell_2}C$ 的元素；
其前已有 $\ell_1$ 个 $Q$ 因子，
故链长应为 $\ell=\ell_1+\ell_2$。
见疑误 A-0551。}
更准确地说，
\[
\vec Q^+
=\sum_{a_1,\ldots,a_{\ell_2}\in A}
\pm\vec Q\otimes
\Delta_{a_1}^{2n}\otimes\cdots
\otimes\Delta_{a_{\ell_2}}^{2n}.
\]

沿用第~\ref{subsec:7-3-2} 小节的记号，
取
$\vec P\in
B_{k-|\vec a|}C(L;\Lambda_{0,\mathrm{nov}})[1]$
和
$\vec P^+\in
B_kC(L;\Lambda_{0,\mathrm{nov}})[1]$
\translatornote{源书此处两个链长分别写成
$k+1-|\vec a|$ 与 $k+1$；
但式 (7.3.33) 的边界评价因子是 $L^k$，
而第 7.3.2 小节的输入链长度分别是
$k-|\vec a|$ 与 $k$。
见疑误 A-0552。}。
考虑纤维积
\begin{equation}
\begin{aligned}
&\mathcal M_{k+1;\ell}(L;\beta)
\times_{M^\ell\times L^k}
(\vec Q^+\times\vec P^+)\\
&\qquad=
\mathcal M_{k+1;\ell}
(L;\beta;\vec Q^+,\vec P^+).
\end{aligned}
\tag{7.3.33}\label{eq:7-3-33}
\end{equation}
要在其上定义由
$[0,1]^{\ell_2}\times[0,1]^{|\vec a|}$
参数化的一族多值截面
$\mathfrak s_{s_1,\ldots,s_{\ell_2};
t_1,\ldots,t_{|\vec a|}}$，
并要求它们满足下列相容条件。

\par\noindent\textup{(7.3.34.1)}\label{cond:7-3-34-1}\quad
当 $t_i\to1$ 或 $t_i\to0$ 时，
相容条件分别与相容条件~\ref{cond:7-3-13}、
\ref{cond:7-3-14}、
\ref{cond:7-3-17} 类似。

\par\noindent\textup{(7.3.34.2)}\label{cond:7-3-34-2}\quad
当 $s_i\to1$ 时，
多值截面收敛到
$\mathcal M_{k+1;\ell}
(L;\beta;\vec Q^{+-},\vec P^+)$
上的多值截面 $\mathfrak s_*$ 的拉回，
其中
$\vec Q^{+-}
=\vec Q\otimes M^{\otimes(\ell_2-1)}$。

\par\noindent\textup{(7.3.34.3)}\label{cond:7-3-34-3}\quad
置 $\vec Q'=\vec Q\otimes M$。
当 $s_i\to0$ 时，
$\mathcal M_{k+1;\ell}
(L;\beta;\vec Q^+,\vec P^+)$
上的多值截面
$\mathfrak s_{s_1,\ldots,s_{\ell_2};
t_1,\ldots,t_{|\vec a|}}$
的拉回收敛到
$\mathcal M_{k+1;\ell}
(L;\beta;\vec Q'{}^+,\vec P^+)$
上的多值截面
$\mathfrak s_{s_1,\ldots,s_{i-1},
s_{i+1},\ldots,s_{\ell_2};
t_1,\ldots,t_{|\vec a|}}$。
这里
$\vec Q'{}^+
=\vec Q'\otimes M^{\otimes(\ell_2-1)}
=\vec Q^+$。

\par\noindent\textup{(7.3.34.4)}\label{cond:7-3-34-4}\quad
在
$\mathcal M_{k+1;\ell}
(L;\beta;\vec Q'{}^+,\vec P^+)$
的边界上，
多值截面应与同类模空间的纤维积相容，
方式与第~\ref{sec:7-2} 节中
其他类似情形相同。

还要求这些多值截面
在内部标记点的置换下不变。
另外还需一项相容条件，
描述内部标记点在极限中
变为边界标记点时多值截面的行为。

取
$((\Sigma,\vec z,\vec z^+),w)
\in\mathcal M_{k+1;\ell}(L;\beta)$。
这里 $\Sigma$ 是亏格为零的有边黎曼曲面，
$\vec z^+=(z_1^+,\ldots,z_\ell^+)$
为内部标记点，
$\vec z=(z_0,z_1,\ldots,z_k)$
为边界标记点，
而
$w\colon(\Sigma,\partial\Sigma)\to(M,L)$
是同调类 $\beta$ 中的伪全纯映射。

把 $z_i\in\partial\Sigma$
与 $1\in\partial D^2$ 识别，
将 $\Sigma$ 与 $D^2$ 粘接。
记所得亏格为零的有边黎曼曲面为
$\Sigma^+$。
它的输入边界标记点少了一个，
余下 $k-1$ 个为
$\vec z\setminus\{z_i\}$
中除输出点 $z_0$ 以外的点；
把其内部 $\ell+1$ 个标记点定义为
$\vec z^+\cup\{0\}$，
其中 $0\in D^2$。
将 $w$ 延拓为
$w^+\colon\Sigma^+\to M$，
使它在新粘上的
$\Sigma^+\setminus\Sigma=D^2$
上恒等于 $w(z_i)$。
于是
\[
((\Sigma,\vec z,\vec z^+),w)
\longmapsto
((\Sigma^+,\vec z\setminus\{z_i\},
\vec z^+\cup\{0\}),w^+)
\]
诱导映射
\[
I_i\colon
\mathcal M_{k+1;\ell}
(L;\beta;\vec Q^+,\vec P^+)
\longrightarrow
\mathcal M_{k;\ell+1}
(L;\beta;\vec Q^{++},\vec P^{+-}),
\]
其中
\[
\begin{aligned}
\vec Q^{++}&=\vec Q^+\otimes M,\\
\vec P^{+-}
&=P_1\otimes\cdots\otimes P_{i-1}
\otimes P_{i+1}\otimes\cdots\otimes P_k.
\end{aligned}
\]
令 $i=\vec a(j)$。
所需的附加相容条件是：
在目标
$\mathcal M_{k;\ell+1}
(L;\beta;\vec Q^{++},\vec P^{+-})$
和源
$\mathcal M_{k+1;\ell}
(L;\beta;\vec Q^+,\vec P^+)$
上分别取相应的多值截面，并要求
\begin{equation}
\begin{aligned}
&I_i^*\!\left(
\mathfrak s_{
\substack{s_1,\ldots,s_{\ell_2+1};\\
t_1,\ldots,t_{j-1},t_{j+1},\ldots,t_{|\vec a|}}}
\right)\\
&\qquad=
\mathfrak s_{
\substack{s_1,\ldots,s_{\ell_2};\\
t_1,\ldots,t_{j-1},s_{\ell_2+1},
t_{j+1},\ldots,t_{|\vec a|}}}.
\end{aligned}
\tag{7.3.34.5}\label{cond:7-3-34-5}
\end{equation}
图~\ref{fig:7-3-1}
示意这一内部标记点变为边界标记点的识别。
\par\medskip\noindent\begin{minipage}{\textwidth}
\centering
\setlength{\unitlength}{0.62mm}
\begin{picture}(200,73)
\linethickness{0.6pt}
{\color{gray!65}
  \qbezier(45,48)(86,48)(86,29)
  \qbezier(86,29)(86,10)(45,10)
  \qbezier(45,10)(4,10)(4,29)
  \qbezier(4,29)(4,48)(45,48)
  \qbezier(151,48)(192,48)(192,29)
  \qbezier(192,29)(192,10)(151,10)
  \qbezier(151,10)(110,10)(110,29)
  \qbezier(110,29)(110,48)(151,48)}
\qbezier(45,43.5)(81,43.5)(81,29)
\qbezier(81,29)(81,14.5)(45,14.5)
\qbezier(45,14.5)(9,14.5)(9,29)
\qbezier(9,29)(9,43.5)(45,43.5)
\qbezier(151,43.5)(187,43.5)(187,29)
\qbezier(187,29)(187,14.5)(151,14.5)
\qbezier(151,14.5)(115,14.5)(115,29)
\qbezier(115,29)(115,43.5)(151,43.5)
\put(45,52){\circle{18}}
\put(45,52){\circle*{1.2}}
\put(45,43.5){\circle*{1.2}}
\put(45,68){$\scriptstyle M$}
\put(45,56){\line(0,1){10}}
\put(96,28){$\Longrightarrow$}
\put(151,43.5){\circle*{1.2}}
\put(148,49){$\scriptstyle L$}
\put(27,26){\circle*{1}}
\put(36,26){\circle*{1}}
\put(45,26){\circle*{1}}
\put(54,26){\circle*{1}}
\put(27,27){\line(0,1){7}}
\put(36,27){\line(0,1){7}}
\put(45,27){\line(0,1){7}}
\put(54,27){\line(0,1){7}}
\put(24,35){$\scriptstyle Q$}
\put(33,35){$\scriptstyle Q$}
\put(42,35){$\scriptstyle Q$}
\put(51,35){$\scriptstyle Q$}
\put(133,26){\circle*{1}}
\put(142,26){\circle*{1}}
\put(151,26){\circle*{1}}
\put(160,26){\circle*{1}}
\put(133,27){\line(0,1){7}}
\put(142,27){\line(0,1){7}}
\put(151,27){\line(0,1){7}}
\put(160,27){\line(0,1){7}}
\put(130,35){$\scriptstyle Q$}
\put(139,35){$\scriptstyle Q$}
\put(148,35){$\scriptstyle Q$}
\put(157,35){$\scriptstyle Q$}
\put(16,47){$\scriptstyle P_{i-1}$}
\put(59,47){$\scriptstyle P_{i+1}$}
\put(122,47){$\scriptstyle P_{i-1}$}
\put(165,47){$\scriptstyle P_{i+1}$}
\put(83,28){$\scriptstyle P_k$}
\put(189,28){$\scriptstyle P_k$}
\put(60,4){$\scriptstyle P_1$}
\put(166,4){$\scriptstyle P_1$}
\put(21,4){$\scriptstyle P$}
\put(127,4){$\scriptstyle P$}
\put(71,13){$\scriptstyle\operatorname{ev}_0$}
\put(177,13){$\scriptstyle\operatorname{ev}_0$}
\end{picture}
\par\smallskip
{\renewcommand{\thefigure}{7.3.1}\refstepcounter{figure}
\normalfont\bfseries 图~\thefigure\label{fig:7-3-1}}
\end{minipage}\par\medskip


与命题~\ref{prop:7-3-10} 同理，
可以构造横截于零、
并满足 (7.3.34.1)--(7.3.34.5)
的一族多值截面
$\mathfrak s_{s_1,\ldots,s_{\ell_2};
t_1,\ldots,t_{|\vec a|}}$。
严格说来，
还须像命题~\ref{prop:7-3-10} 那样
选择光滑奇异单形的可数集等；
这里略去细节。

现在定义
\[
\begin{aligned}
&\mathfrak q^{++}_{\ell,k+1}
\bigl(
\vec Q\otimes
\widetilde{\mathbf f}^{\otimes\ell_2},
P_1\otimes\cdots\otimes\mathbf f
\otimes\cdots\otimes\mathbf f
\otimes\cdots\otimes P_k
\bigr)\\
&\qquad=
\sum_\beta T^{\omega(\beta)}
e^{\mu_L(\beta)/2}
\operatorname{ev}_{0*}
\left[
\mathcal M
(L;\beta;\vec Q^+,\vec P^+)^{\mathfrak s}
\right].
\end{aligned}
\]
把 (3.8.36.1)--(3.8.36.3)
作为 $\mathfrak q^{++}$ 的定义的一部分。
它满足式~\eqref{eq:3-8-33}，
由 (7.3.34.1)--(7.3.34.4) 得出；
式~\eqref{eq:3-8-36-4}
则由 (7.3.34.5) 得出
\translatornote{源书此处写 (3.8.36.5)，
但定理 3.8.32 的该组公式
只有 (3.8.36.1)--(3.8.36.4)；
这里应指末项 (3.8.36.4)。
见疑误 A-0553。}。
其细节与命题~\ref{prop:7-3-10}
的证明类似，故从略。
至此完成 (3.8.36) 的证明。
